\abstract{%
Modern sequence models---including genomic language models (GLMs) such as Enformer, Borzoi, and Evo~2, as well as long-context large language models---accept nominal context windows spanning hundreds of kilobases to megabases (or hundreds of thousands of tokens), yet empirical analyses repeatedly indicate that only a fraction of the input context is functionally utilized for any given prediction. This discrepancy motivates a rigorous, estimable, and comparable notion of \emph{effective context length} (ECL): a model- and data-dependent quantity measuring the input span that materially affects a model's embedding or output at a reference locus.

We formalize ECL for embedding-generating models as a property of the joint distribution of sequences and perturbations. Our primary definition---\emph{perturbation--variance ECL}---is the minimal radius around a reference position needed to capture a prescribed fraction of the expected squared embedding change induced by localized, distribution-aware perturbations. We introduce several companion quantities: the \emph{Effective Context Profile} (ECP), the \emph{Effective Context Dimension} (ECD), \emph{directional ECL}, \emph{interaction ECL}, and \emph{Context Decay Spectroscopy} (CDS). We present a unified theoretical treatment linking perturbation ECL to variance-based global sensitivity indices (Sobol/total-effect), effective receptive field theory, information-theoretic context measures, and decision-theoretic model comparison via Blackwell sufficiency. We establish monotonicity, invariance properties, exponential decay bounds, minimax estimation rates, and architectural locality results. We develop statistically consistent estimators with non-asymptotic concentration guarantees (Bernstein bounds), antithetic variance reduction, importance sampling, and bootstrap confidence intervals. We propose a comprehensive experimental protocol---including a genomic Needle-in-a-Haystack test (gNIAH)---for estimating and comparing ECL across genomic models and tasks. Eight appendices provide complete proofs, implementation guidance, computational complexity analysis, and an annotated reference guide.
}
